\documentclass[aps,pra,reprint,amsmath,amssymb,nofootinbib,floatfix]{revtex4-2}
\usepackage[T1]{fontenc}
\usepackage{graphicx}
\usepackage{bm}
\usepackage[colorlinks=true,citecolor=blue,linkcolor=blue,urlcolor=blue]{hyperref}
\newcommand{\Tr}{\operatorname{Tr}}
\newcommand{\Rea}{\operatorname{Re}}
\newcommand{\Ima}{\operatorname{Im}}
\newcommand{\dd}{\mathrm{d}}
\newcommand{\ii}{\mathrm{i}}
\newcommand{\ee}{\mathrm{e}}
\newcommand{\ket}[1]{\lvert #1\rangle}
\newcommand{\bra}[1]{\langle #1\rvert}
\newtheorem{proposition}{Proposition}
\newenvironment{proof}{\par\noindent\textit{Proof. }}{\hfill$\square$\par}

\begin{document}
\title{Weak periodic phase shaping lowers the worst-case sensitivity of coherent-state qubit readout}
\author{Xiaojian Zeng}
\email{xiaojianzeng@foxmail.com}
\affiliation{PBC School of Finance, Tsinghua University, 100190, China}
\date{October 12, 2026}

\begin{abstract}
A fixed interaction can encode the real amplitude of a coherent bosonic
mode into a qubit, with the field subsequently discarded. The optimal
linear quadrature coupling gives amplitude-independent quantum Fisher
information $4/\ee$. We show that every fixed real bounded periodic phase
perturbation along the same qubit rotation axis, nonconstant almost
everywhere, lowers the infimum of the symmetric-logarithmic-derivative
quantum Fisher information over all real amplitudes, for both sufficiently
small nonzero signs of its strength. A Gaussian-score calculation gives a
nonzero, zero-mean first-order sensitivity response whose Fourier
multipliers have no nonzero-frequency zeros. Explicit uniform rank and
quadratic-remainder bounds convert this response into a strict worst-case
loss. The argument includes discontinuous profiles and requires neither a
photon cutoff nor a sampled-amplitude extrapolation. The strength radius
depends on the fixed profile and period. This restricted obstruction
excludes a weak periodic shaping route; it does not determine the
unrestricted interaction optimum or exclude finite-strength improvements.
\end{abstract}
\maketitle

\section{Introduction}

When a quantized field is measured through a small probe, the interaction
must transfer information into the degrees of freedom that are actually
read out. For the real amplitude $\alpha$ of a single coherent mode, the
field quantum Fisher information (QFI) is $4$, whereas a fixed qubit
readout generally retains less. Ravell Rodr\'{i}guez, Perarnau-Llobet, and
Sekatski formulated the optimization of the worst-amplitude qubit QFI and
exhibited a linear quadrature coupling with uniform value
$4/\ee$~\cite{ravell2026}. Their work leaves the general fixed-interaction
optimization unresolved and suggests $4/\ee$ or a nearby value.

A natural design question is whether weak nonlinear shaping of that
linear interaction can improve its uniform sensitivity. Optimizing the
linear coupling strength alone does not answer this question: nonlinear
profiles introduce infinitely many additional directions. Here we analyze
each fixed bounded periodic profile acting along the same qubit rotation
axis. Every nonconstant profile produces a first-order ripple in the local
QFI, and that ripple necessarily has both signs. A uniform remainder
estimate then proves a strict reduction of the worst-amplitude QFI for
either sufficiently small sign of the perturbation.

This is a directionwise statement. The profile and its period are fixed
before the perturbation strength tends to zero. The result does not
provide a common neighborhood over all profiles, an upper bound for
finite-strength shaping, or a solution of the arbitrary-unitary problem.

\section{Model and the linear benchmark}

The field Hilbert space is $\ell^2(\mathbb N_0)$, with $[a,a^\dagger]=1$.
In oscillator units, $\hbar=1$ and
\begin{equation}
 Q=\frac{a+a^\dagger}{\sqrt2},\qquad
 p_\alpha(x)=\frac{\ee^{-(x-\sqrt2\alpha)^2}}{\sqrt\pi}.
 \label{eq:gaussian}
\end{equation}
The input is the real coherent state $\ket\alpha$, so its quadrature has
mean $\sqrt2\alpha$ and variance $1/2$. The amplitude, quadrature, and
couplings below are dimensionless. The qubit begins in the fixed state
$\ket g$. In the basis $(\ket g,\ket e)$, use
$\sigma_y=\left(\begin{smallmatrix}0&-\ii\\\ii&0\end{smallmatrix}\right)$,
so $\ee^{-\ii\sigma_y f(x)}\ket g=
\cos f(x)\ket g+\sin f(x)\ket e$.

Let $L>0$ and let $g$ be a real bounded Borel $L$-periodic function, fixed
independently of $\alpha$ and of the strength $\epsilon$. We consider
\begin{equation}
 V_\epsilon=\exp\left[-\ii\sigma_y\otimes
 \left(\frac{Q}{\sqrt2}+\epsilon g(Q)\right)\right].
 \label{eq:unitary}
\end{equation}
Since $g(Q)$ is bounded and self-adjoint, the quadrature function in
Eq.~\eqref{eq:unitary} is self-adjoint on the domain of $Q$.
Functional calculus defines a unitary on the full field space, including
for discontinuous $g$. Only the final qubit is retained:
\begin{align}
 \rho_\epsilon(\alpha)
 &=\Tr_F\left[V_\epsilon
 (\ket g\bra g\otimes\ket\alpha\bra\alpha)
 V_\epsilon^\dagger\right]\notag\\
 &=\int_{\mathbb R}p_\alpha(x)
 \ket{\psi_\epsilon(x)}\bra{\psi_\epsilon(x)}\,\dd x,
 \label{eq:state}\\
 \ket{\psi_\epsilon(x)}
 &=\ee^{-\ii\sigma_y[x/\sqrt2+\epsilon g(x)]}\ket g.
 \notag
\end{align}
There is one input mode per encoding, no postselection or retained
classical branch record, and no control adapted to the unknown amplitude.

We use the pointwise symmetric logarithmic derivative (SLD) QFI
\cite{braunstein1994}, defined by
\begin{equation}
 \rho_\epsilon\mathcal L_\alpha+
 \mathcal L_\alpha\rho_\epsilon=2\partial_\alpha\rho_\epsilon,
 \qquad F_Q=\Tr(\rho_\epsilon\mathcal L_\alpha^2).
 \label{eq:sld}
\end{equation}
This is a local sensitivity benchmark, allowing optimization of the final
qubit measurement at a local working point. It does not imply a globally
identifiable amplitude or a particular finite-sample estimation risk.
Uniform full rank, established below, avoids rank-changing SLD conventions.

Put $\mu=\sqrt2\alpha$ and $p_0(y)=\ee^{-y^2}/\sqrt\pi$. The complex
Bloch coordinate is
\begin{equation}
 m_\epsilon(\mu)=\int_{\mathbb R}p_0(y)
 \ee^{\ii\sqrt2(\mu+y)+2\ii\epsilon g(\mu+y)}\,\dd y.
 \label{eq:m}
\end{equation}
In particular,
\begin{equation}
 \rho_\epsilon=\frac12
 \begin{pmatrix}1+\Rea m_\epsilon&\Ima m_\epsilon\\
 \Ima m_\epsilon&1-\Rea m_\epsilon\end{pmatrix}.
 \label{eq:blochstate}
\end{equation}
Define $d_\epsilon(\mu)=\sqrt2\partial_\mu m_\epsilon(\mu)$ and
$\mathcal F_\epsilon(\mu)=F_Q[\rho_\epsilon(\mu/\sqrt2)]$. Thus
$d_\epsilon$ is always an \emph{amplitude} derivative. Solving
Eq.~\eqref{eq:sld} in the Bloch representation gives
\begin{equation}
 \mathcal F_\epsilon(\mu)=|d_\epsilon|^2+
 \frac{[\Rea(\overline m_\epsilon d_\epsilon)]^2}
 {1-|m_\epsilon|^2}.
 \label{eq:qfi}
\end{equation}
For completeness, writing $\mathcal L_\alpha=b_0 I+\bm b\cdot\bm\sigma$
gives $b_0+\bm b\cdot\bm r=0$ and
$\bm b+b_0\bm r=\partial_\alpha\bm r$, which immediately yields
Eq.~\eqref{eq:qfi} for $\bm r=(\Ima m,0,\Rea m)$.

More generally, the linear coupling $f(x)=\tau x/\sqrt2$ gives
\begin{equation}
 m_\tau(\alpha)=\ee^{-\tau^2/2+2\ii\tau\alpha},
 \qquad F_Q=4\tau^2\ee^{-\tau^2}.
 \label{eq:baseline}
\end{equation}
This is the exact quadrature benchmark of Ref.~\cite{ravell2026},
not a large-amplitude approximation to the Jaynes--Cummings interaction.
It is maximized at $|\tau|=1$, giving $4/\ee$ for every real $\alpha$.
Equation~\eqref{eq:unitary} perturbs the $\tau=1$ carrier.

\section{The first-order sensitivity response}

The derivative in Eq.~\eqref{eq:qfi} can be obtained by differentiating
the Gaussian density at fixed quadrature $x$:
\begin{equation}
 d_\epsilon(\mu)=\int_{\mathbb R}2\sqrt2\,y\,p_0(y)
 \ee^{\ii\sqrt2(\mu+y)+2\ii\epsilon g(\mu+y)}\,\dd y.
 \label{eq:score}
\end{equation}
The integrable Gaussian score justifies this differentiation even when
$g$ has jumps; no derivative of $g$ is used.

Let $A=\ee^{-1/2}$. At zero strength,
\begin{equation}
 m_0=A\ee^{\ii\sqrt2\mu},\qquad
 d_0=2\ii A\ee^{\ii\sqrt2\mu}.
 \label{eq:zerodata}
\end{equation}
Hence $|d_0|^2=4/\ee$ and
$\Rea(\overline m_0d_0)=0$. If $m_1=\partial_\epsilon m_\epsilon|_0$
and $d_1=\partial_\epsilon d_\epsilon|_0$, then
\begin{align}
 m_1&=2\ii\int p_0(y)g(\mu+y)
 \ee^{\ii\sqrt2(\mu+y)}\,\dd y,\notag\\
 d_1&=4\ii\sqrt2\int y p_0(y)g(\mu+y)
 \ee^{\ii\sqrt2(\mu+y)}\,\dd y.
 \label{eq:firstdata}
\end{align}
The radial term in Eq.~\eqref{eq:qfi} starts at second order. Therefore
the first variation is
\begin{align}
 G(\mu)&=\left.\partial_\epsilon\mathcal F_\epsilon(\mu)\right|_0
 =2\Rea(\overline d_0d_1)\notag\\
 &=16\sqrt2 A\int_{\mathbb R}y p_0(y)
 \cos(\sqrt2y)g(\mu+y)\,\dd y.
 \label{eq:response}
\end{align}
The kernel is smooth, integrable, and odd. Its convolution with bounded
$g$ is continuous, $L$-periodic, and has zero mean.

To determine whether any nonconstant profile can have zero response,
write $\kappa=2\pi/L$ and
\begin{equation}
 c_n=\frac1L\int_0^L g(x)\ee^{-\ii n\kappa x}\,\dd x.
\end{equation}
The Gaussian moment
$\int y p_0(y)\ee^{\ii\lambda y}\,\dd y
=\ii\lambda\ee^{-\lambda^2/4}/2$ gives, for $n>0$,
\begin{equation}
 G_n=\ii C(n\kappa)c_n,\qquad
 G_{-n}=-\ii C(n\kappa)c_{-n},\qquad G_0=0,
 \label{eq:fourierresponse}
\end{equation}
where $k_0=1/\sqrt2$ and
\begin{align}
 C(\xi)&=16k_0\ee^{-1-\xi^2/4}\notag\\
 &\quad\times\left[\xi\cosh(k_0\xi)-2k_0\sinh(k_0\xi)\right].
 \label{eq:multiplier}
\end{align}
Fubini's theorem applies because $g$ is bounded and the kernel is
integrable. No pointwise Fourier-series convergence of $g$ is assumed.

For $\xi>0$, put $z=k_0\xi$. The bracket in
Eq.~\eqref{eq:multiplier} is
$\sqrt2(z\cosh z-\sinh z)$, which is strictly positive: it vanishes
at $z=0$ and has derivative $\sqrt2 z\sinh z>0$.
Thus $C(\xi)>0$ at every nonzero positive frequency.
If $G$ vanished identically, Eq.~\eqref{eq:fourierresponse} would force
all nonzero Fourier coefficients of $g$ to vanish. Fourier uniqueness
in $L^2([0,L])$ would then make $g$ constant almost everywhere.
For a nonconstant profile, it follows that
\begin{align}
 \min_\mu G(\mu)&<0<\max_\mu G(\mu),\notag\\
 \gamma&=\min\{-\min G,\max G\}>0.
 \label{eq:gamma}
\end{align}
The extrema exist on any compact period. Each fixed perturbation therefore
has an amplitude at which its first-order QFI change is unfavorable,
for either sign of strength.

For a finite trigonometric profile,
\begin{equation}
g(x)=g_0+\sum_{n\geq1}[a_n\cos(n\kappa x)+b_n\sin(n\kappa x)],
\end{equation}
the response is explicitly
\begin{equation}
 G(\mu)=\sum_{n\geq1}C(n\kappa)
 [b_n\cos(n\kappa\mu)-a_n\sin(n\kappa\mu)].
 \label{eq:trigresponse}
\end{equation}
Distinct harmonics cannot cancel the entire ripple: no nonzero mode lies
in the response kernel.

\section{Uniform rank control and strict worst-case loss}

A first variation at isolated amplitudes would not by itself control the
whole real axis. We next bound the QFI remainder uniformly. Choose any
finite $M>0$ satisfying $|g(x)|\leq M$, and define
\begin{align}
 D&=2A,\qquad W=2\sqrt{2/\pi},\qquad h=(1-A)/4,\notag\\
 \beta&=(1+A)/2,\qquad \delta=1-\beta^2,\notag\\
 K&=4W^2(1+h)^2+4DW+
 \frac{4(\beta W+D)^2}{\delta}.
 \label{eq:constants}
\end{align}
All these constants are positive, and $\beta<1$.

Write $a=|\epsilon|$. The inequalities
$|\ee^{2\ii\epsilon g}-1|\leq2aM$ and
$\int 2\sqrt2|y|p_0(y)\,\dd y=W$ imply
\begin{equation}
 |m_\epsilon-m_0|\leq2aM,\qquad
 |d_\epsilon-d_0|\leq2aMW.
 \label{eq:lipschitz}
\end{equation}
When $aM\leq h$, these give
\begin{equation}
 |m_\epsilon|\leq A+2h=\beta,
 \qquad 1-|m_\epsilon|^2\geq\delta.
 \label{eq:rank}
\end{equation}
The smallest output eigenvalue is at least $(1-\beta)/2=h$,
independently of the amplitude.

For real $g$, the exponential Taylor remainder obeys
$|\ee^{2\ii\epsilon g}-1-2\ii\epsilon g|\leq2a^2M^2$.
Consequently,
\begin{align}
 m_\epsilon&=m_0+\epsilon m_1+r_m,&|r_m|&\leq2a^2M^2,\notag\\
 d_\epsilon&=d_0+\epsilon d_1+r_d,&|r_d|&\leq2a^2M^2W,\notag\\
 |m_1|&\leq2M,&|d_1|&\leq2MW.
 \label{eq:taylor}
\end{align}
Using $|d_0|=D$ and $G=2\Rea(\overline d_0d_1)$, expansion
of the first term of Eq.~\eqref{eq:qfi} yields
\begin{align}
 \big||d_\epsilon|^2-D^2-\epsilon G\big|
 &\leq2D|r_d|+|\epsilon d_1+r_d|^2\notag\\
 &\leq a^2M^2[4DW+4W^2(1+h)^2].
 \label{eq:tangentialbound}
\end{align}
For the second term, the zero radial derivative at the carrier and
Eq.~\eqref{eq:lipschitz} give
\begin{align}
 |\Rea(\overline m_\epsilon d_\epsilon)|
 &\leq |m_\epsilon|\,|d_\epsilon-d_0|
       +|m_\epsilon-m_0|\,|d_0|\notag\\
 &\leq2aM(\beta W+D).
 \label{eq:radialbound}
\end{align}
Combining Eqs.~\eqref{eq:rank}, \eqref{eq:tangentialbound}, and
\eqref{eq:radialbound} proves the exact bound
\begin{equation}
 \left|\mathcal F_\epsilon(\mu)-\frac4\ee-\epsilon G(\mu)\right|
 \leq KM^2\epsilon^2
 \quad\text{for all }\mu\in\mathbb R,
 \label{eq:uniformremainder}
\end{equation}
whenever $|\epsilon|M\leq h$.

\begin{proposition}
For each fixed $L>0$ and real bounded Borel $L$-periodic profile $g$
that is not constant almost everywhere, let $\gamma$ and $K$ be given
by Eqs.~\eqref{eq:gamma} and \eqref{eq:constants}. With
\begin{equation}
 \epsilon_0(g,L)=\min\left\{\frac hM,
                  \frac{\gamma}{2KM^2}\right\}>0,
 \label{eq:radius}
\end{equation}
every $0<|\epsilon|<\epsilon_0$, of either sign, gives uniformly
full-rank qubit outputs and
\begin{equation}
 \inf_{\alpha\in\mathbb R}F_Q[\rho_\epsilon(\alpha)]<\frac4\ee.
 \label{eq:mainresult}
\end{equation}
\end{proposition}

\begin{proof}
For $\epsilon>0$, evaluate Eq.~\eqref{eq:uniformremainder} at a minimum
of $G$. For $\epsilon<0$, use a maximum of $G$. In either case, using
$\mu=\sqrt2\alpha$ and Eq.~\eqref{eq:radius},
\begin{align}
 \inf_{\alpha\in\mathbb R}F_Q[\rho_\epsilon(\alpha)]
 &\leq\frac4\ee-\gamma|\epsilon|+KM^2\epsilon^2\notag\\
 &<\frac4\ee-\frac\gamma2|\epsilon|<\frac4\ee.
 \label{eq:lossbound}
\end{align}
Equation~\eqref{eq:rank} supplies the uniform rank claim.
\end{proof}

If $g$ is constant almost everywhere, the perturbation is a fixed output
rotation and preserves $4/\ee$ exactly. Thus the exclusion of constant
profiles is necessary, and changes on sets of measure zero are immaterial.

\section{Illustration and numerical consistency checks}

For $g(x)=\sin(\kappa x)$, Eq.~\eqref{eq:trigresponse} reduces to
$G(\mu)=C(\kappa)\cos(\kappa\mu)$. Figure~\ref{fig:ripple} illustrates
the response and the opposite unfavorable amplitude regions for the two
signs. Its QFI values were calculated from the Gaussian projector and
score integrals, followed by a spectral SLD evaluation. They are supporting
numerical data; the proposition follows from the uniform analytic bound.
The illustrated strengths satisfy the rank/remainder condition
$|\epsilon|M\leq h$; they need not lie within the conservative strict-loss
radius in Eq.~\eqref{eq:radius}.

\begin{figure}[!htbp]
 \includegraphics[width=\columnwidth]{manuscript/figures/sensitivity-ripple.pdf}
 \caption{Sensitivity ripple for the fixed profile
 $g(x)=\sin(\kappa x)$ with $\kappa=0.7$, plotted against
 $t=\kappa\mu=\sqrt2\kappa\alpha$.
 QFI change from $4/\ee$ at $\epsilon=0.005$ and $\epsilon=-0.02$,
 from direct density-matrix SLD calculations at 17 phases. Lines join sampled values
 to guide the eye. Both curves have values below the baseline. No grid
 minimum or numerical error enclosure is used to prove the all-amplitude
 statement.}
 \label{fig:ripple}
\end{figure}

The same consistency tests used a mixed-harmonic profile, a discontinuous
square wave, a high harmonic, and a constant control. Across five profiles
and 17 phases at two strengths, the 170 direct SLD evaluations satisfy
the analytic rank and remainder inequalities. The largest discrepancy
between the Gaussian and Fourier first-response calculations is below
$3.8\times10^{-15}$, and the largest observed ratio of the absolute QFI
remainder to $KM^2\epsilon^2$ is below $0.057$. The square-wave integral
is split at its jumps. These finite floating-point checks do not provide
outward-rounded quadrature certificates or establish the response's
all-frequency injectivity. Equations~\eqref{eq:multiplier}--\eqref{eq:uniformremainder}
provide the corresponding exact statements. Code, the saved data, and
regeneration instructions accompany this draft.

\section{Physical consequence and limitations}

The loss mechanism is a nonzero first-order modulation with zero average:
weak shaping can increase the QFI at some amplitudes, but decreases it at
others. The response filter is injective on nonconstant periodic profiles,
so adding harmonics cannot eliminate the entire first-order modulation.
This rules out every fixed weak bounded periodic same-axis profile,
including profiles with jumps, as a route to improving the established
uniform qubit readout.

The result concerns one fixed interaction over the whole amplitude axis.
Locally optimized preprocessing under noisy measurements, such as the
finite-dimensional setting of Ref.~\cite{zhou2023}, has a different
objective and does not furnish this uniform fixed-profile conclusion.
Likewise, retaining an additional measurement record or querying several
field modes would change the readout resources specified here.

The radius in Eq.~\eqref{eq:radius} depends on $g$ and $L$.
The multipliers can be small at low or high frequencies, so the present
argument supplies no common radius over all periods or profiles.
Finite-strength deformations, strength-dependent profiles, and interactions
outside the same-axis functional form remain beyond the proposition.
Their possible success is not established. The unrestricted fixed-unitary
worst-amplitude optimization is not resolved by this result, and no
experimental implementation or global estimation advantage is claimed.

\section*{Use of AI tools}

OpenAI Codex was used substantively for literature retrieval, derivation
development, symbolic and numerical calculations, code preparation,
manuscript drafting, and internal checking. Separate AI sessions
reconstructed the source model and recomputed the restricted result.
These checks are not external peer review or human scientific verification.
Human scientific review and authorship confirmation remain pending for
this local draft; the accompanying author handoff records the disclosure
details that require confirmation.

\section*{Data and code availability}

The accompanying reproducibility material contains the analytic argument,
numerical consistency-test code and saved outputs, figure-generation
script, software versions, and build instructions. No experimental data
were used.

\bibliographystyle{apsrev4-2}
\bibliography{manuscript/references}
\end{document}
